\section*{अध्याय \ref{ch:b2:cell-signalling} --- रासायनिक संदेशवाहक और संकेत-पारक्रमण}

\begin{solution}{exo:b2:cell-signalling:1}
इंसुलिन: अंतःस्रावी, जल-घुलनशील (पेप्टाइड)। कॉर्टिसोल: अंतःस्रावी,
वसा-घुलनशील। किसी सिनैप्स पर ऐसिटिलकोलीन: सिनैप्टिक, जल-घुलनशील।
किसी घाव पर हिस्टामिन: पराक्राइन, जल-घुलनशील। ईस्ट्राडायोल:
अंतःस्रावी, वसा-घुलनशील। अधिवृक्क से ऐड्रिनैलिन: अंतःस्रावी,
जल-घुलनशील। नाइट्रिक ऑक्साइड: पराक्राइन, वसा-घुलनशील (यानी ऐसी गैस जो
झिल्लियाँ पार कर जाती है)।
\end{solution}

\begin{solution}{exo:b2:cell-signalling:2}
ऐड्रिनैलिन उस ग्राही से बँधता है (बंद: वियोजन, ग्राही का
भीतर खींचा जाना); ग्राही \omterm{prop:b2:cell-signalling:gpcr}{G प्रोटीन} सक्रिय करता है, GDP के बदले GTP (बंद:
सेकंडों में GTP का जल-अपघटन); \omterm{prop:b2:cell-signalling:gpcr}{G प्रोटीन} ऐडिनिलिल साइक्लेज़ सक्रिय करता है,
जो cAMP बनाता है (बंद: फ़ॉस्फ़ोडाइएस्टरेज़); cAMP काइनेज़ A सक्रिय करता है
(बंद: cAMP की हानि); काइनेज़ A फ़ॉस्फ़ोरिलेज़ काइनेज़ का फ़ॉस्फ़ोरिलन करता है,
जो ग्लाइकोजन फ़ॉस्फ़ोरिलेज़ का फ़ॉस्फ़ोरिलन करता है (बंद: फ़ॉस्फ़ेटेज़);
फ़ॉस्फ़ोरिलेज़ ग्लाइकोजन से ग्लूकोज़-1-फ़ॉस्फ़ेट छोड़ता है, जो बदलकर
ग्लूकोज़ के रूप में बाहर भेजी जाती है।
\end{solution}

\begin{solution}{exo:b2:cell-signalling:3}
सतही ग्राही जल-घुलनशील संदेशवाहकों को झिल्ली पर बाँधते हैं और
\omterm{prop:b2:cell-signalling:gpcr}{द्वितीय संदेशवाहकों} तथा काइनेज़ों से होकर सेकंडों में काम करते हैं, और
मौजूदा प्रोटीनों की क्रियाशीलता बदलते हैं। \omterm{prop:b2:cell-signalling:nuclear}{केंद्रकीय ग्राही} वसा-घुलनशील
संदेशवाहकों को कोशिका के भीतर बाँधते हैं और \omterm{prop:b2:cell-differentiation:transcriptional}{अनुलेखन कारकों} के रूप में घंटों में
काम करते हैं, और बदलते हैं कि कौन-से प्रोटीन बनें।
\end{solution}

\begin{solution}{exo:b2:cell-signalling:4}
\omterm{prop:b2:cell-signalling:gpcr}{चक्रीय AMP}: ऐडिनिलिल साइक्लेज़ बनाता है,
फ़ॉस्फ़ोडाइएस्टरेज़ नष्ट करता है। कैल्सियम: वाहिकाओं से घुसता है या
इनोसिटॉल ट्राइफ़ॉस्फ़ेट के उत्तर में अंतर्द्रव्यी जालिका छोड़ता है, और
पंप हटाते हैं। इनोसिटॉल ट्राइफ़ॉस्फ़ेट (डाइएसिलग्लिसरॉल के साथ): किसी
झिल्ली-लिपिड से फ़ॉस्फ़ोलिपेज़ C बनाता है, और फ़ॉस्फ़ेटेज़ हटाते हैं।
\end{solution}

\begin{solution}{exo:b2:cell-signalling:5}
$\theta = [L]/(K_d + [L])$: $0.09$, $0.5$, $0.91$, $0.99$। ग्राहियों को
दुगुना करना वह अंश अपरिवर्तित छोड़ता है और अधिग्रहीत संख्या ---
और इस तरह वह अनुक्रिया --- दुगुनी कर देता है।
\end{solution}

\begin{solution}{exo:b2:cell-signalling:6}
$50\times 200\times 100\times 300 = 3\times 10^{8}$। सौ अधिग्रहीत
ग्राही प्रति सेकंड $3\times 10^{10}$ उत्पाद-अणु देते हैं --- यानी कोई
ऊपरी सीमा, क्योंकि वे चरण संतृप्त हो जाते हैं।
\end{solution}

\begin{solution}{exo:b2:cell-signalling:7}
वह जकड़ा \omterm{prop:b2:cell-signalling:gpcr}{G प्रोटीन} साइक्लेज़ चालू रखता है, cAMP ऊँची बनी रहती है, काइनेज़ A
आँत-कोशिकाओं की क्लोराइड वाहिका फ़ॉस्फ़ोरिलित तथा
खुली रखता है, क्लोराइड स्रवित होता है, और सोडियम तथा जल उसके पीछे
गुहा में जाते हैं --- प्रतिदिन लीटरों भर। वह जकड़ \omterm{prop:b2:cell-signalling:gpcr}{G प्रोटीन} का कोई
सहसंयोजी रूपांतर है, जो केवल तब उलटता है जब वे कोशिकाएँ स्वयं बदल दी जाएँ, कुछ दिन
बाद।
\end{solution}

\begin{solution}{exo:b2:cell-signalling:8}
$0.9\times 10^{-6}\,\text{mol/L}\times 2\times 10^{-12}\,\text{L} =
1.8\times 10^{-18}$ मोल, यानी $1.1\times 10^{6}$ आयन (असल में और भी, क्योंकि
बफ़र जो घुसे उसका अधिकांश बाँध लेते हैं)। प्रति सेकंड $10^{6}$ पर एक वाहिका
को सेकंड भर लगता; और हज़ार वाहिकाएँ यह मिलीसेकंड भर में कर देती हैं,
और वही काल-मान किसी संकेत को चाहिए।
\end{solution}

\begin{solution}{exo:b2:cell-signalling:9}
वे दोनों पथ उन्हीं एंजाइमों पर आ मिलते हैं: ग्लूकागॉन का काइनेज़ A
ग्लाइकोजन सिंथेज़ (बंद) तथा फ़ॉस्फ़ोरिलेज़ काइनेज़ (चालू) का
फ़ॉस्फ़ोरिलन करता है; इंसुलिन का सोपान वह फ़ॉस्फ़ेटेज़ सक्रिय करता है जो वे
फ़ॉस्फ़ेट हटाता है। उन एंजाइमों की दशा काइनेज़ तथा
फ़ॉस्फ़ेटेज़ क्रियाशीलता के संतुलन से बैठती है, इसलिए जो गिनता है वह दोनों
\omterm{def:b2:cell-signalling:kinds}{हॉर्मोनों} का अनुपात है; दोनों ऊँचे हों तो आम तौर पर वह फ़ॉस्फ़ेटेज़ जीतता है और ग्लाइकोजन
संचित होता है।
\end{solution}

\begin{solution}{exo:b2:cell-signalling:10}
ऐड्रिनैलिन: cAMP \qty{20}{\micro m} की किसी कोशिका के आरपार सेकंड भर से
कम में विसरित हो जाता है ($x^{2}/2D$, जिसमें $D \approx \qty{3e-10}{m^2/s}$), और उस
सोपान का हर एंजाइम पहले से मौजूद है --- यानी सेकंड। कॉर्टिसोल: कोई
अनुलेख \qty{60}{s} लेता है, कोई प्रोटीन \qty{100}{s}; और कुछ
दर्जन mRNA से, जिनमें से हर एक दर्जन भर राइबोसोम ढो रहा हो, वह कोशिका प्रति मिनट कुछ
सौ एंजाइम-अणु बनाती है, इसलिए किसी मापन-योग्य प्रभाव को चाहिए हज़ारों में घंटे लगते
हैं, और वह भी उन मिनटों के बाद जो कॉर्टिसोल को घुसने तथा DNA तक पहुँचने में लगते हैं।
\end{solution}

\begin{solution}{exo:b2:cell-signalling:11}
वह ग्राही लगातार संकेत देता है: MAP काइनेज़ सोपान उस
कोशिका को बिना किसी वृद्धि कारक के विभाजित होने को चलाता है, और वह उस
संकेत की अनुपस्थिति की उपेक्षा करता है जो सामान्यतः उसे रोके रखता है। चूँकि कोई एक
\omterm{def:b2:genome-diversification:mutation}{उत्परिवर्तन} विभाजन पर से कोई नियंत्रण हटा देता है, इसलिए ऐसे ग्राही (और नीचे के
काइनेज़) सबसे आम अर्बुदजनों में हैं। ऐसी कोई दवा जो
उस काइनेज़ के ATP स्थल पर बैठ जाए वे फ़ॉस्फ़ोरिलन रोक देती है और
उस ग्राही को चुप कर देती है --- यही कई कैंसर-दवाओं का सिद्धांत है।
\end{solution}

\begin{solution}{exo:b2:cell-signalling:12}
अंतःस्रावी: मिनट भर में हर कोशिका तक पहुँचता है, प्रति कोशिका सस्ता, किसी एक
कोशिका को संबोधित नहीं कर सकता, मिनटों से दिनों तक काम करता है --- यानी
पूरी-देह की दशाओं (उपवास, वृद्धि, प्रजनन) के लिए अपरिहार्य। तंत्रिका:
मिलीसेकंड, एक लक्ष्य, महँगी तार-जोड़, मिलीसेकंडों तक काम करती है
--- यानी गति तथा तेज़ प्रतिवर्तों के लिए अपरिहार्य। दोनों उसी संदेश के लिए
वहाँ काम में लिए जाते हैं जहाँ गति तथा पहुँच दोनों चाहिए: अनुकंपी
तंत्रिकाएँ तथा अधिवृक्क मध्यांश दोनों किसी डर में
नॉरऐड्रिनैलिन या ऐड्रिनैलिन पहुँचाते हैं।
\end{solution}

\begin{solution}{pb:b2:cell-signalling:1}
\textbf{1.} $1\,\text{nmol}/5\,\text{L} = \qty{0.2}{nmol/L}$।
\textbf{2.} $\theta = 0.2/1.2 = 0.17$: यानी लगभग 3300 अधिग्रहीत ग्राही।
\textbf{3.} तीन अर्ध-आयु: \qty{0.025}{nmol/L}; $\theta = 0.024$।
\textbf{4.} लगभग आधे मिनट से मिनट भर, यानी जितना समय \omterm{def:b2:blood-circulation:blood}{रक्त}
घूमने में लेता है; जबकि वह तंत्रिका अपना संचारी सीधे
यकृत तक सेकंड भर में पहुँचा देती है।
\textbf{5.} वही अंश, पर $33\,000$ अधिग्रहीत ग्राही: यानी
दस गुना बड़ी अनुक्रिया।
\textbf{6.} \qty{1}{nmol/L} पर $K_d = \qty{1}{\micro
mol/L}$ वाला कोई ग्राही हज़ारवाँ भाग अधिग्रहीत है: कोई अनुक्रिया नहीं। $K_d$ को
परिसंचरित परास से मिलाना अधिग्रहण-वक्र का तीखा भाग वहाँ रखता है जहाँ
वह \omterm{def:b2:cell-signalling:kinds}{हॉर्मोन} सचमुच बदलता है।
\textbf{7.} $20\times 100\times 1\times 50\times 50\times 1000 =
5\times 10^{9}$ ग्लूकोज़ अणु प्रति सेकंड प्रति अधिग्रहीत ग्राही।
\textbf{8.} पहले सेकंड में $3300\times 20\times 100 = 6.6\times 10^{6}$
cAMP; \qty{5e-12}{L} में, $1.1\times 10^{-17}$ मोल, यानी लगभग
\qty{2}{\micro mol/L}।
\textbf{9.} $3300\times 5\times 10^{9} = 1.7\times 10^{13}$ प्रति सेकंड, यानी प्रति कोशिका $10^{15}$
प्रति मिनट।
\textbf{10.} $2\times 10^{11}$ कोशिकाएँ $\times 10^{15} = 2\times
10^{26}$ अणु प्रति मिनट, यानी 330 मोल, यानी प्रति मिनट \qty{6e4}{g} --- यानी
देखे गए \qty{5}{g/min} से दस हज़ार गुना। वे लब्धियाँ केवल तभी गुणा होती हैं
जब हर चरण असंतृप्त हो; असल में वह फ़ॉस्फ़ोरिलेज़
अपने क्रियाधार तथा फ़ॉस्फ़ेटेज़ों से सीमित है, और ग्लूकोज़ का बाहर भेजा जाना
वाहकों से। उस सोपान की लब्धि गति तथा
संवेदिता पर ख़र्च होती है, निर्गम पर नहीं।
\textbf{11.} \qty{5}{g/min} पर बीस मिनट; वह डर मिनटों में ख़त्म हो जाता है
और वे बंद-स्विच उस सोपान को उससे बहुत पहले रोक देते हैं।
\textbf{12.} \omterm{prop:b2:cell-signalling:gpcr}{G प्रोटीन} द्वारा GTP का जल-अपघटन (सेकंडों में);
cAMP पर फ़ॉस्फ़ोडाइएस्टरेज़ (सेकंडों में); फ़ॉस्फ़ोरिलित एंजाइमों पर
फ़ॉस्फ़ेटेज़ (सेकंडों से मिनटों में);
ग्राही-असंवेदीकरण (मिनटों में)। सबसे तेज़ \omterm{prop:b2:cell-signalling:gpcr}{G प्रोटीन} की अपनी घड़ी
और वह फ़ॉस्फ़ोडाइएस्टरेज़ हैं।
\textbf{13.} cAMP नष्ट नहीं होता: वह जमा होता और बना रहता है, वह
काइनेज़ चालू रहता है, और वह ग्लूकोज़-निर्गम बड़ा होता तथा अधिक देर चलता है
--- यानी कॉफ़ी की वह सतर्कता।
\textbf{14.} $4000/50 = \qty{80}{s}$।
\textbf{15.} प्रति राइबोसोम प्रति एंजाइम $600/5 = \qty{120}{s}$; यानी प्रति
राइबोसोम प्रति घंटा 30, और प्रति mRNA प्रति घंटा 600।
\textbf{16.} प्रति घंटा $100\times 600 = 60\,000$ एंजाइम; और $10^{5}$
लगभग \qty{1.7}{h} बाद।
\textbf{17.} कॉर्टिसोल को घुसने, बँधने और अपने जीनों तक पहुँचने में दस से बीस
मिनट, पहले mRNA के बहुसंख्य तथा अनूदित होने से पहले घंटा भर,
फिर दो घंटे का जमाव: यानी लगभग चार घंटे।
\textbf{18.} कोई सोपान उन एंजाइमों की क्रियाशीलता बदलता है जो मौजूद हैं;
जबकि कॉर्टिसोल का काम यह बदलना है कि कौन-से एंजाइम मौजूद हों, और वह भी दिनों के लिए, जो
केवल अनुलेखन कर सकता है। कोई डर नए एंजाइमों के लिए चार घंटे प्रतीक्षा नहीं
कर सकता।
\textbf{19.} अधिक ग्राहियों का अर्थ है किसी भी ऐड्रिनैलिन-सांद्रता पर
अधिक अधिग्रहीत ग्राही: वह अधिग्रहण-अंश अपरिवर्तित है पर
हर सांद्रता पर वह अनुक्रिया बड़ी है --- यानी वह वक्र
ऊपर पैमाना बदलता है, खिसकता नहीं। कोई धीमा \omterm{def:b2:cell-signalling:kinds}{हॉर्मोन} किसी तेज़ की लब्धि बाँध देता है।
\textbf{20.} ग्लूकोज़ $\beta$ कोशिका में घुसती है और उसका उपापचय होता है; ATP
चढ़ती है और कोई ATP-संवेदी पोटैशियम वाहिका बंद कर देती है; वह झिल्ली
विध्रुवित होती है; वोल्टता-द्वारित कैल्सियम वाहिकाएँ खुलती हैं; कैल्सियम
इंसुलिन-कणिकाओं का बहिःकोशिकन छेड़ती है --- यानी किसी विद्युत् चरण से
स्रवण से युग्मित उपापचय।
\textbf{21.} वे दोनों ग्राही दो पथ शुरू करते हैं जो उन्हीं एंजाइमों पर
आ मिलते हैं: ग्लूकागॉन का काइनेज़ A उनका फ़ॉस्फ़ोरिलन करता है, इंसुलिन का
सोपान वह फ़ॉस्फ़ेटेज़ सक्रिय करता है जो उनका विफ़ॉस्फ़ोरिलन करता है; और उन
एंजाइमों की दशा, तथा इस तरह ग्लाइकोजन-उपापचय की दिशा,
दोनों के अनुपात के पीछे चलती है।
\textbf{22.} उपवास-ग्लूकोज़ ऊँची (यकृत, बिना रोक-टोक, उसे
बनाता रहता है), भोजन के बाद ग्लूकोज़ बहुत ऊँची (पेशी तथा वसा द्वारा कोई ग्रहण नहीं);
और इंसुलिन के बिना वसा-भंडार वसीय अम्ल छोड़ते हैं, जिन्हें
यकृत कीटोन अम्लों में बदल देता है --- यानी कीटोअम्लरक्तता।
\textbf{23.} काम करती पेशी इंसुलिन के बिना ग्लूकोज़ ग्रहण करती है, इसलिए किसी
लंबी दौड़ पर वह ग्लूकोज़ गिरती है जबकि इंजेक्ट की गई इंसुलिन, बिना नियमन,
यकृत को दबाती रहती है: यानी अल्पग्लूकोज़रक्तता तथा ढह जाना।
\textbf{24.} ग्लूकोज़-पाश: संवेदक $\beta$ तथा $\alpha$ कोशिकाएँ,
दोनों दिशाओं के लिए दो प्रभावक \omterm{def:b2:cell-signalling:kinds}{हॉर्मोन}, कोई ऊँची लब्धि जो उस स्तर को
पाँचवें भाग के भीतर थामे, और काम करने में मिनट। \omterm{def:b2:blood-pressure:baroreflex}{दाब-प्रतिवर्त}:
खिंचाव-ग्राही, प्रभावक के रूप में \omterm{def:b2:heart:anatomy}{हृदय} तथा वाहिकाएँ, लगभग चार की लब्धि, और काम करने में
सेकंड भर।
\textbf{25.} छोड़े जाने के बाद अधिग्रहण $0.17$; प्रति सेकंड प्रति अधिग्रहीत ग्राही
लब्धि $5\times 10^{9}$; असल में यकृत-निर्गम \qty{5}{g/min}; और कॉर्टिसोल
लगभग चार घंटे लेता है।
\end{solution}
